\documentclass[10pt,a4paper]{amsart}
\usepackage{amsmath,amssymb,amsthm,graphicx}
% For commutative diagrams we use Xy-pic
\usepackage[all]{xy}

% We use 2cell for 2-commutative diagrams.
\xyoption{2cell}
% Theorem environments.
%
\theoremstyle{plain}
\newtheorem{theorem}[subsection]{Theorem}
\newtheorem{proposition}[subsection]{Proposition}
\newtheorem{lemma}[subsection]{Lemma}

\theoremstyle{definition}
\newtheorem{definition}[subsection]{Definition}
\newtheorem{example}[subsection]{Example}
\newtheorem{exercise}[subsection]{Exercise}
\newtheorem{situation}[subsection]{Situation}

\theoremstyle{remark}
\newtheorem{remark}[subsection]{Remark}
\newtheorem{remarks}[subsection]{Remarks}

\numberwithin{equation}{subsection}

% Macros
%
\def\lim{\mathop{\mathrm{lim}}\nolimits}
\def\colim{\mathop{\mathrm{colim}}\nolimits}
\def\Spec{\mathop{\mathrm{Spec}}}
\def\Hom{\mathop{\mathrm{Hom}}\nolimits}
\def\Ext{\mathop{\mathrm{Ext}}\nolimits}
\def\SheafHom{\mathop{\mathcal{H}\!\mathit{om}}\nolimits}
\def\SheafExt{\mathop{\mathcal{E}\!\mathit{xt}}\nolimits}
\def\Sch{\mathit{Sch}}
\def\Mor{\mathop{\mathrm{Mor}}\nolimits}
\def\Ob{\mathop{\mathrm{Ob}}\nolimits}
\def\Sh{\mathop{\mathit{Sh}}\nolimits}
\def\NL{\mathop{N\!L}\nolimits}
\def\CH{\mathop{\mathrm{CH}}\nolimits}
\def\proetale{{pro\text{-}\acute{e}tale}}
\def\etale{{\acute{e}tale}}
\def\QCoh{\mathit{QCoh}}
\def\Ker{\mathop{\mathrm{Ker}}}
\def\Im{\mathop{\mathrm{Im}}}
\def\Coker{\mathop{\mathrm{Coker}}}
\def\Coim{\mathop{\mathrm{Coim}}}

% Boxtimes
%
\DeclareMathSymbol{\boxtimes}{\mathbin}{AMSa}{"02}

%
% Macros for moduli stacks/spaces
%
\def\QCohstack{\mathcal{QC}\!\mathit{oh}}
\def\Cohstack{\mathcal{C}\!\mathit{oh}}
\def\Spacesstack{\mathcal{S}\!\mathit{paces}}
\def\Quotfunctor{\mathrm{Quot}}
\def\Hilbfunctor{\mathrm{Hilb}}
\def\Curvesstack{\mathcal{C}\!\mathit{urves}}
\def\Polarizedstack{\mathcal{P}\!\mathit{olarized}}
\def\Complexesstack{\mathcal{C}\!\mathit{omplexes}}
% \Pic is the operator that assigns to X its picard group, usage \Pic(X)
% \Picardstack_{X/B} denotes the Picard stack of X over B
% \Picardfunctor_{X/B} denotes the Picard functor of X over B
\def\Pic{\mathop{\mathrm{Pic}}\nolimits}
\def\Picardstack{\mathcal{P}\!\mathit{ic}}
\def\Picardfunctor{\mathrm{Pic}}
\def\Deformationcategory{\mathcal{D}\!\mathit{ef}}
\usepackage{cleveref}
\crefname{originallemma}{original lemma}{original lemmas}
\Crefname{originallemma}{Original lemma}{Original lemmas}
\makeatletter
\newlabel{lemma-describe-pullback}{{09XN}{source}{}{sourcepointer-1}{}}
\newlabel{lemma-describe-pullback@cref}{{[originallemma][0][]09XN}{[1][0][]source}{}{}{}}
\makeatother
\title{Etale sections over the base agree with sections over a flat quasi-compact cover with connected geometric fibres}
\author{The Stacks Project Authors}
\date{stacks-project:etale-cohomology.tex:local-snapshot-20261011; GNU FDL 1.2 or later}
\begin{document}
\maketitle
\noindent\textit{Editorial note.} All delivered mathematical text is the author's own, copied from the pinned Stacks Project chapter file etale-cohomology.tex; no statement or proof step was rewritten, reordered or removed.\par\medskip
\noindent\textit{Editorial note.} The proof is the author's complete argument: the injection, the fpqc reduction and the geometric-fibre computation are all inside the counted span, and the author's closing phrase is retained verbatim.\par\medskip
\noindent\textit{Editorial note.} The label lemma-sections-upstairs occurs in no other staged Stacks chapter, so no cross-chapter restatement had to be weighed for this card.\par\medskip
\noindent\textit{Editorial note.} The single labelled reference occurring inside the counted spans is carried in assembly.external\_references as a non-load-bearing navigation pointer to the same chapter.\par\medskip
\noindent\textit{Editorial note.} The preamble is the chapter's own native preamble excerpt already used for this source: the Xy-pic setup (native preamble lines 14-18) and the native theorem declarations with the mathematical macros (native preamble lines 154-217).\par\medskip
\noindent\textit{Allowed standard prerequisites (editorial).}\begin{itemize}
\item Small inverse image f\_\{small\}\textasciicircum{}\{-1\} of a sheaf on the small etale site. Same chapter, the pullback f\_\{small\}\textasciicircum{}\{-1\} of a sheaf on S\_etale along a morphism f : T -> S. Supplies the sheaf whose sections over T are compared with the sections over S.
\item fpqc coverings and the description of the pullback of a section. Same chapter, the fact that a flat quasi-compact surjective morphism gives an fpqc covering and the resulting description of pullbacks of sections. Supplies the reduction of surjectivity to agreement of the two pullbacks to T x\_S T.
\item Constant sheaf on a connected scheme and constancy of its sections. Same chapter, the constant sheaf with value F\_sbar on the small etale site of a connected scheme, whose sections are constant. Supplies the identification of the two pullbacks over each geometric fibre, which is the crux of the proof.
\end{itemize}\medskip
\ifdefined\hypertarget\hypertarget{sourcepointer-1}{}\fi
\noindent\textit{Source reference.} Original lemma 09XN: sources/stacks-etale-cohomology/pinned-original-source.tex, source lines 4960--4960. Same-chapter navigation pointer invoked at proof lines 5049-5050, where the author uses the description of the pullback of a section along an fpqc covering. The pointer only resolves that reference for the reader; the cited lemma is not reproduced inside the counted span, and no bibliography citation is involved.\par\medskip
\begin{lemma}
\label{lemma-sections-upstairs}
Let $S$ be a scheme. Let $f : T \to S$ be a morphism such that
\begin{enumerate}
\item $f$ is flat and quasi-compact, and
\item the geometric fibres of $f$ are connected.
\end{enumerate}
Let $\mathcal{F}$ be a sheaf on $S_\etale$.
Then $\Gamma(S, \mathcal{F}) = \Gamma(T, f^{-1}_{small}\mathcal{F})$.
\end{lemma}

\begin{proof}
There is a canonical map
$\Gamma(S, \mathcal{F}) \to \Gamma(T, f_{small}^{-1}\mathcal{F})$.
Since $f$ is surjective (because its fibres are connected) we see that
this map is injective.

\medskip\noindent
To show that the map is surjective, let
$\alpha \in \Gamma(T, f_{small}^{-1}\mathcal{F})$.
Since $\{T \to S\}$ is an fpqc covering we can use
Lemma \ref{lemma-describe-pullback} to see that suffices to prove that
$\alpha$ pulls back to the same section over $T \times_S T$ by the
two projections. Let $\overline{s} \to S$ be a geometric point.
It suffices to show the agreement holds over $(T \times_S T)_{\overline{s}}$
as every geometric point of $T \times_S T$ is contained in one of
these geometric fibres. In other words, we are trying to show that
$\alpha|_{T_{\overline{s}}}$ pulls back to the same section over
$$
(T \times_S T)_{\overline{s}} =
T_{\overline{s}} \times_{\overline{s}} T_{\overline{s}}
$$
by the two projections to $T_{\overline{s}}$.
However, since $\mathcal{F}|_{T_{\overline{s}}}$ is the
pullback of $\mathcal{F}|_{\overline{s}}$ it is a constant sheaf
with value $\mathcal{F}_{\overline{s}}$. Since $T_{\overline{s}}$
is connected by assumption, any section of a constant sheaf is constant.
Hence $\alpha|_{T_{\overline{s}}}$ corresponds to an element
of $\mathcal{F}_{\overline{s}}$. Thus the two pullbacks to
$(T \times_S T)_{\overline{s}}$ both
correspond to this same element and we conclude.
\end{proof}

\end{document}
